\documentclass{article} %report, book, .... article
\usepackage[utf8]{inputenc} %\ddddddd[]{}
\usepackage{kotex}
\usepackage{color}
\usepackage{amssymb}
\usepackage{amsmath}
\usepackage{mathtools}
\usepackage{extarrows}
\usepackage{graphicx}
\usepackage[hypcap=false]{caption}
\usepackage{indentfirst}
\usepackage{array}
\usepackage{enumitem}
\usepackage{multirow}
\usepackage{makecell}
\usepackage{fancyhdr}
\usepackage{appendix}
\usepackage{mathrsfs}
\usepackage{listings}
\usepackage{algorithm}
\usepackage{algpseudocode}
\newcolumntype{L}[1]{>{\raggedright\let\newline\\\arraybackslash\hspace{0pt}}m{#1}}
\newcolumntype{C}[1]{>{\centering\let\newline\\\arraybackslash\hspace{0pt}}m{#1}}
\newcolumntype{R}[1]{>{\raggedleft\let\newline\\\arraybackslash\hspace{0pt}}m{#1}}
\usepackage[a4paper,left=25mm,right=25mm,top=30mm,bottom=30mm]{geometry}
\usepackage{xcolor}
\usepackage[normalem]{ulem} % provide \uline without changing \emph
\usepackage{hyperref}
\usepackage{pgfplots}
\RequirePackage{tikz}
\RequirePackage[framemethod=TikZ]{mdframed}

\pgfplotsset{compat=1.18}
\definecolor{mplblue}{RGB}{31,119,180} % 원본 느낌을 살리는 Matplotlib 기본 파란색
% Global: don't change text color or draw borders for links so TOC and refs look normal
\hypersetup{
	colorlinks=false,
	pdfborder={0 0 0}
}
% Redefine inline link commands to show blue text with a blue underline.
% We capture the original commands then wrap the link text.
\let\OriginalHref\href
\renewcommand{\href}[2]{% #1 = URL, #2 = text
	\OriginalHref{#1}{\textcolor{blue}{\uline{#2}}}%
}

\let\OriginalUrl\url
\renewcommand{\url}[1]{% typeset a clickable, blue-underlined URL
	\href{#1}{\textcolor{blue}{\uline{\nolinkurl{#1}}}}%
}
\usepackage{multicol}

%\pagestyle{myheadings}
\pagestyle{fancy}
\fancyhf{}
\fancyhead[R]{(2026S) Term Project - Proposal}
\fancyhead[L]{CS.20200 - Programming Prinsiples}
% \fancyhead[R]{}
\fancyfoot[L]{KAIST}
\fancyfoot[C]{\thepage}
\fancyfoot[R]{20240131 Seoho Kim}
\title{\huge\textbf{2048 CLI}}
\author{20240131 Seoho Kim}

\linespread{1.6}
\setcounter{tocdepth}{2}

% {\"a} {\^e} {\`i} {\.I} {\o} {\'u} {\aa} {\c c} {\u g} {\l} {\~n} {\H o} {\v r} {\s s} {\r u}

\setlist[enumerate, 1]{label=\textbf{R\arabic*.}, leftmargin=*, itemsep=4pt}
% \setlist[enumerate, 1]{label=(\alph*)}
% \setlist[enumerate, 2]{label=\roman*.}

\usepackage{amsthm}
\newtheorem*{theorem}{Theorem}
\newtheorem*{definition}{Definition}
\newtheorem*{lemma}{Lemma}
\newtheorem*{claim}{\textbf{Claim}}

% ---------- listings (for inline code style) ----------

\lstset{
    basicstyle=\ttfamily\small,
    backgroundcolor=\color{gray!10},
    frame=single,
    framerule=0pt,
    breaklines=true
}

\definecolor{plGray}{RGB}{180, 180, 180}

\global\mdfdefinestyle{defaultStyle}{
        topline=false, bottomline=false, rightline=false,
        leftline=true,
        linewidth=3pt,
        backgroundcolor=plBackground,
        frametitleaboveskip=10pt,
        frametitlebelowskip=0pt,
        frametitlealignment=\raggedright,
        % frametitlefont=\bfseries\sffamily,
        innerleftmargin=12pt, innerrightmargin=12pt,
        innertopmargin=5pt, innerbottommargin=12pt,
        skipabove=15pt, skipbelow=15pt,
        nobreak=true,
}

\newenvironment{myNote}{
        \begin{mdframed}[
                        style=defaultStyle,
                        linecolor=plGray,
                        backgroundcolor=plGray!8,
                        innertopmargin=7pt,
                        frametitle={\textbf{\color{plGray}Note}}
                ]
                \setlength{\parindent}{0pt}
                \setlength{\parskip}{5pt}
                }{
        \end{mdframed}
}

% ================================================================
\begin{document}
\maketitle
\thispagestyle{fancy}

% ----------------------------------------------------------------
\section{Overview}

This project is a command-line implementation of the sliding-tile puzzle game \textbf{2048},
written entirely in \textbf{F\# with .NET~10}.
The player slides tiles on a square board using the \texttt{W}/\texttt{A}/\texttt{S}/\texttt{D} keys.
When two tiles of the same value collide, they merge into a single tile whose value is their sum.
The player wins when a tile reaches the value~2048.
The game is lost when no legal move remains.

The player may choose from four board sizes before each game:
$2{\times}2$, $4{\times}4$, $8{\times}8$, and $16{\times}16$.
The per-size all-time high score is persisted to a local file between sessions.

The implementation follows a FUNCTIONAL CORE and IMPERATIVE SHELL architecture:
all game logic (tile merging, move validation, win/loss detection) is expressed as
pure, immutable functions with no mutable state, while console I/O is confined to the
outermost layer.


% ----------------------------------------------------------------
\section{Requirements}

\subsection{Main Menu}

\begin{enumerate}

	\item When the program starts, it displays a main menu containing exactly four options:\\
	      \texttt{1) 3x3}, \texttt{2) 4x4}, \texttt{3) 5x5}, \texttt{4) 6x6}, \texttt{5) 8x8}.

	\item The user selects a board size by entering the digit
	      \texttt{1}, \texttt{2}, \texttt{3}, \texttt{4}, or \texttt{5}.
	      If any other input is entered, the menu is re-displayed and the user is prompted again.

	\item The main menu also displays the all-time high score for each board size,
	      loaded from the local score file.
	      If no score has been recorded for a size, it is displayed as~\texttt{0}.

\end{enumerate}

\subsection{Game Initialisation}

\begin{enumerate}[resume]

	\item After the user selects a board size,
	      a new game board of that size is created with all cells initialised to~\texttt{0}.

	\item Exactly two tiles are placed on the board in randomly chosen distinct empty cells
	      before the first move.
	      Each tile's value is independently sampled:
	      \texttt{2} with probability~$0.9$ and \texttt{4} with probability~$0.1$.

	\item The current score is initialised to~\texttt{0}.

\end{enumerate}

\subsection{Display}

\begin{enumerate}[resume]

	\item After every state change (initialisation, successful move, win, or loss),
	      the board is re-rendered to the terminal.

	\item Each cell is displayed in a consistent format such that all cells have equal width.
	      The exact formatting, padding, and representation of empty cells are implementation-defined.

	\item The current score and the all-time high score for the selected board size
	      are displayed on every render.

	\item The available controls (\texttt{W}/\texttt{A}/\texttt{S}/\texttt{D} to move,
	      \texttt{R} to restart, \texttt{Q} to quit) are displayed on every render.

\end{enumerate}

\subsection{Moves}

\begin{enumerate}[resume]

	\item During a game the program reads exactly one character of input at a time without
	      requiring the user to press Enter (raw keypress mode).
	      Valid move keys are \texttt{W} (up), \texttt{A} (left), \texttt{S} (down), \texttt{D} (right),
	      case-insensitively.

	\item A \emph{slide} in a given direction is computed row-by-row (for left/right) or
	      column-by-column (for up/down).
	      Within each line, the following steps are applied in order:
	      \begin{enumerate}[label=(\alph*)]
		      \item Remove all zero cells, compacting non-zero values toward the leading edge.
		      \item Scan the compacted sequence from the leading edge: whenever two adjacent
		            values are equal, replace them with a single cell whose value is their sum.
		            Each cell participates in at most one merge per move.
		      \item Pad the resulting sequence with zeros at the trailing edge to restore the
		            original line length.
	      \end{enumerate}

	\item When two tiles merge, their sum is added to the current score.

	\item If the resulting board after a slide is identical to the board before the slide,
	      the move is considered \emph{ineffective}.
	      An ineffective move does not place a new tile and does not change the score.

	\item If a move is effective, one new tile is placed in a randomly chosen empty cell.
	      The tile value is \texttt{2} with probability~$0.9$ and \texttt{4} with probability~$0.1$.

\end{enumerate}

\subsection{Win Condition}

\begin{enumerate}[resume]

	\item If any cell on the board contains the value~\texttt{2048} after a move, the game
	      enters the \emph{win} state.

	\item Upon entering the win state, the board and score are rendered, followed by the message
	      \texttt{You win!}

	\item The current score is compared with the stored high score for the selected board size.
	      If the current score is greater, the high score is updated in the local score file.

	\item After displaying the win message, the program waits for any keypress,
	      then returns to the main menu.

\end{enumerate}

\subsection{Loss Condition}

\begin{enumerate}[resume]

	\item After every effective move and tile placement, the program checks whether any legal
	      move remains.
	      A legal move exists if and only if at least one of the four slides would produce a board
	      different from the current board.

	\item If no legal move exists, the game enters the \emph{loss} state.

	\item Upon entering the loss state, the board and score are rendered, followed by the message
	      \texttt{Game over.}

	\item The high score update rule in R17 applies identically on loss.

	\item After displaying the game-over message, the program waits for any keypress,
	      then returns to the main menu.

\end{enumerate}

\subsection{In-Game Controls}

\begin{enumerate}[resume]

	\item Pressing \texttt{R} (case-insensitively) during a game immediately ends the current game
	      without updating the high score, and returns the player to the main menu.

	\item Pressing \texttt{Q} (case-insensitively) during a game immediately ends the current game
	      without updating the high score, and exits the program.

	\item Any key other than \texttt{W}, \texttt{A}, \texttt{S}, \texttt{D}, \texttt{R}, \texttt{Q}
	      is silently ignored; the board is not re-rendered.

\end{enumerate}

\subsection{Score Persistence}

\begin{enumerate}[resume, label=\textbf{R\arabic*.}, leftmargin=*, itemsep=4pt]

	\item High scores are stored in a plain-text file named \texttt{scores.txt} in the same
	      directory as the executable.

	\item The file contains one line per board size in the format \texttt{<size>:<score>},
	      e.g.\ \texttt{4:13204}.
	      Board sizes not yet recorded are omitted; their high score is treated as~\texttt{0}.

	\item If \texttt{scores.txt} does not exist at program startup, it is treated as empty
	      (all high scores are~\texttt{0}).
	      The file is created when the first high score is written.

\end{enumerate}

% ----------------------------------------------------------------
\section{Example Interaction}

\begin{myNote}
	The following interaction is illustrative. The exact visual layout may differ in the final implementation.
\end{myNote}

The program starts and displays:

\begin{lstlisting}
=== 2048-CLI ===
Select board size:
  1) 2x2   (Best: 0)
  2) 4x4   (Best: 0)
  3) 8x8   (Best: 0)
  4) 16x16 (Best: 0)
> _
\end{lstlisting}

The user presses \texttt{2}.
A $4{\times}4$ board is created, two tiles are placed, and the following is rendered:

\begin{lstlisting}
Score: 0   Best: 0
.  .  .  2
.  .  .  .
.  4  .  .
.  .  .  .
[W/A/S/D] Move  [R] Restart  [Q] Quit
\end{lstlisting}

The user presses \texttt{A} (slide left).
Tiles move left; a new tile appears in a random empty cell.
The updated board and score are displayed.

% ----------------------------------------------------------------
\section{Out of Scope}

The following features are explicitly \emph{not} part of this project's requirements
and will not be implemented:

\begin{itemize}[itemsep=2pt]
	\item Graphical or browser-based rendering.
	\item Online or networked leaderboards.
	\item Undo functionality.
	\item Animations or sound effects.
	\item Targets other than~2048 (e.g., a ``continue past 2048'' mode).
\end{itemize}

% ----------------------------------------------------------------
\section{Repository}

\url{https://github.com/seohokim-hoya/2048-cli}

\end{document}
